\documentclass{article}
\usepackage{amsmath,amssymb,amsthm}
\usepackage{algorithm}
\usepackage{algorithmic}
\usepackage{enumitem}
\usepackage{hyperref}
\newtheorem{theorem}{Theorem}
\newtheorem{lemma}{Lemma}
\newtheorem{assumption}{Assumption}
\newcommand{\EE}{\mathbb{E}}
\newcommand{\RR}{\mathbb{R}}

\begin{document}

\section*{Cyclic Stochastic Gradient Descent (Cyclic‑SGD)}

\section*{Mathematical Formulation}
Let $f:\RR^{d}\rightarrow\RR$ be the (possibly non‑convex) objective and let 
$\{z_{t}\}_{t\ge0}$ denote an i.i.d. data stream.  
At iteration $t$ we draw a stochastic gradient
\[
g_{t}\;:=\;\nabla f(w_{t};z_{t})\in\RR^{d}.
\]
The learning rate follows a deterministic cyclic schedule of period $2M$:
\[
\eta_{t}\;=\;\eta_{\min}
   +(\eta_{\max}-\eta_{\min})\;
      \Bigl(1-\Bigl|\frac{(t\bmod 2M)}{M}-1\Bigr|\Bigr),
   \qquad t=0,1,\dots
\tag{1}
\]
i.e. a triangular wave that linearly ramps from $\eta_{\min}$ to $\eta_{\max}$
in the first $M$ steps and back to $\eta_{\min}$ in the next $M$ steps.
The update rule is
\[
w_{t+1}=w_{t}-\eta_{t}\,g_{t},\qquad t=0,1,2,\dots
\tag{2}
\]

\section*{Pseudocode}
\begin{algorithm}[H]
\caption{Cyclic‑SGD}
\begin{algorithmic}[1]
\STATE \textbf{Input:} $w_{0}\in\RR^{d}$, $\eta_{\min}>0$, $\eta_{\max}\ge\eta_{\min}$, cycle length $M\in\mathbb N$,
max. iterations $T$.
\FOR{$t=0$ \TO $T-1$}
    \STATE Sample $z_{t}$ and compute stochastic gradient $g_{t}=\nabla f(w_{t};z_{t})$.
    \STATE Compute cyclic step size
    \[
    \eta_{t}\leftarrow \eta_{\min}+(\eta_{\max}-\eta_{\min})
                \Bigl(1-\Bigl|\frac{(t\bmod 2M)}{M}-1\Bigr|\Bigr).
    \]
    \STATE Update $w_{t+1}=w_{t}-\eta_{t} g_{t}$.
\ENDFOR
\STATE \textbf{Output:} $w_{T}$ (or optionally the averaged iterate $\bar w_T=\frac1T\sum_{t=0}^{T-1}w_{t}$).
\end{algorithmic}
\end{algorithm}

\section*{Dry Run Example}
Consider the one–dimensional quadratic
\[
f(w)=(w-3)^{2},\qquad \nabla f(w)=2(w-3).
\]
Set $w_{0}=0$, $\eta_{\min}=0.05$, $\eta_{\max}=0.1$, $M=2$ (so the schedule over the first three steps is
$\eta_{0}=0.05$, $\eta_{1}=0.075$, $\eta_{2}=0.1$).  
Assume the stochastic gradient is exact (i.e. $g_{t}=\nabla f(w_{t})$).

\begin{center}
\begin{tabular}{c|c|c|c}
$t$ & $\eta_{t}$ & $g_{t}=2(w_{t}-3)$ & $w_{t+1}=w_{t}-\eta_{t}g_{t}$ \\ \hline
0 & $0.05$ & $2(0-3)=-6$ & $0-0.05(-6)=0.30$ \\
1 & $0.075$ & $2(0.30-3)=-5.40$ & $0.30-0.075(-5.40)=0.705$ \\
2 & $0.10$ & $2(0.705-3)=-4.590$ & $0.705-0.10(-4.590)=1.164$ \\
\end{tabular}
\end{center}
Corresponding objective values:
$f(w_{0})=9$, $f(w_{1})\approx7.62$, $f(w_{2})\approx3.43$, $f(w_{3})\approx3.01$.
The example illustrates how the cyclic step sizes influence the magnitude of each update.

\newpage

\section*{Assumptions}
\begin{assumption}[Smoothness]\label{ass:smooth}
$f$ is $L$‑smooth: for all $x,y\in\RR^{d}$
\[
f(y)\le f(x)+\langle \nabla f(x),y-x\rangle+\frac{L}{2}\|y-x\|^{2}.
\tag{A1}
\]
\end{assumption}
\begin{assumption}[Unbiased Stochastic Gradient]\label{ass:unbiased}
For every $t$, conditioned on $w_{t}$,
\[
\EE[g_{t}\mid w_{t}]=\nabla f(w_{t}).
\tag{A2}
\]
\end{assumption}
\begin{assumption}[Bounded Variance]\label{ass:variance}
There exists $\sigma^{2}>0$ such that
\[
\EE\big[\|g_{t}-\nabla f(w_{t})\|^{2}\mid w_{t}\big]\le\sigma^{2},
\qquad\forall t.
\tag{A3}
\]
\end{assumption}
\begin{assumption}[Cyclic Step‑size Bounds]\label{ass:step}
\[
0<\eta_{\min}\le \eta_{t}\le\eta_{\max}<\frac{1}{L},
\qquad\forall t.
\tag{A4}
\]
\end{assumption}
All constants $L,\sigma,\eta_{\min},\eta_{\max}$ are known a priori.

\section*{Main Convergence Theorem}
\begin{theorem}[Average Gradient Norm for Cyclic‑SGD]\label{thm:main}
Let $\{w_{t}\}_{t=0}^{T}$ be generated by Cyclic‑SGD under Assumptions \ref{ass:smooth}–\ref{ass:step}.
Define the average step size
\[
\bar\eta:=\frac{1}{T}\sum_{t=0}^{T-1}\eta_{t}.
\]
Then
\[
\frac{1}{T}\sum_{t=0}^{T-1}\EE\big[\|\nabla f(w_{t})\|^{2}\big]
\;\le\;
\frac{2\big(f(w_{0})-f^{\star}\big)}{\bar\eta\,T}
\;+\;
L\big(\eta_{\max}^{2}+\eta_{\min}^{2}\big)\frac{\sigma^{2}}{\bar\eta},
\tag{3}
\]
where $f^{\star}=\inf_{w}f(w)$. In particular, if $\eta_{\max}=O(T^{-1/2})$ and $\eta_{\min}=O(T^{-1/2})$, then
\[
\frac{1}{T}\sum_{t=0}^{T-1}\EE\big[\|\nabla f(w_{t})\|^{2}\big]=O\!\left(\frac{1}{\sqrt{T}}\right).
\]
\end{theorem}

\section*{Theoretical Properties}
\begin{itemize}[leftmargin=*]
\item \textbf{Stability.} Because $\eta_{t}<1/L$, the quadratic term in the smoothness inequality never dominates, guaranteeing that the objective does not increase in expectation.
\item \textbf{Hyper‑parameter Sensitivity.} The bound (3) depends on the *average* step size $\bar\eta$ and on the *range* $\eta_{\max}^{2}+\eta_{\min}^{2}$. Wider cycles (large $M$) keep $\bar\eta$ close to $(\eta_{\min}+\eta_{\max})/2$, while extreme amplitudes increase the variance term.
\item \textbf{Comparison to Classical SGD.} If we set $\eta_{t}\equiv\eta$ (i.e. $M=1$ and $\eta_{\min}=\eta_{\max}=\eta$) then (3) reduces to the standard SGD bound $\frac{2(f(w_{0})-f^{\star})}{\eta T}+L\eta\sigma^{2}$.
\item \textbf{Special Cases.}
    \begin{enumerate}
        \item \emph{Constant schedule:} $\eta_{\min}=\eta_{\max}=\eta$ recovers the classical rate.
        \item \emph{Vanishing amplitude:} $\eta_{\max}\downarrow\eta_{\min}$ yields the same bound with $\eta_{\max}^{2}+\eta_{\min}^{2}\approx2\eta^{2}$.
    \end{enumerate}
\end{itemize}

\section*{Discussion}
The theorem shows that a deterministic cyclic learning‑rate does not deteriorate the $O(1/\sqrt{T})$ rate typical for non‑convex SGD, provided the amplitude of the cycle shrinks at the usual $T^{-1/2}$ speed. The analysis is essentially identical to that of standard SGD; the only extra work is to keep track of the two‑point bound $\eta_{\min}\le\eta_{t}\le\eta_{\max}$ when summing over a period.

\newpage

\section*{Preliminaries (Definitions and Lemmas)}
\begin{lemma}[Smoothness Descent Inequality]\label{lem:smooth}
If $f$ is $L$‑smooth, then for any $w$ and any vector $d$,
\[
f(w-d)\le f(w)-\langle\nabla f(w),d\rangle+\frac{L}{2}\|d\|^{2}.
\tag{L1}
\]
\end{lemma}
\begin{proof}
Directly from Definition \ref{ass:smooth} with $y=w-d$.
\end{proof}

\section*{Main Proof (Step‑by‑Step Derivation)}
We prove Theorem \ref{thm:main}. Throughout, let $\mathcal{F}_{t}$ denote the sigma‑algebra generated by $\{w_{0},\dots,w_{t}\}$.

\noindent\textbf{[Step 1]} Apply Lemma \ref{lem:smooth} with $d=\eta_{t}g_{t}$:
\[
f(w_{t+1})\le f(w_{t})-\eta_{t}\langle\nabla f(w_{t}),g_{t}\rangle
               +\frac{L\eta_{t}^{2}}{2}\|g_{t}\|^{2}.
\tag{4}
\]

\noindent\textbf{[Step 2]} Take conditional expectation w.r.t.\ $\mathcal{F}_{t}$.
Using unbiasedness (A2):
\[
\EE\!\big[f(w_{t+1})\mid\mathcal{F}_{t}\big]
\le f(w_{t})-\eta_{t}\|\nabla f(w_{t})\|^{2}
   +\frac{L\eta_{t}^{2}}{2}\EE\!\big[\|g_{t}\|^{2}\mid\mathcal{F}_{t}\big].
\tag{5}
\]

\noindent\textbf{[Step 3]} Decompose $\|g_{t}\|^{2}$:
\[
\|g_{t}\|^{2}
 =\|\nabla f(w_{t})\|^{2}
   +\|g_{t}-\nabla f(w_{t})\|^{2}
   +2\langle\nabla f(w_{t}),g_{t}-\nabla f(w_{t})\rangle.
\]
Taking conditional expectation and using $\EE[g_{t}-\nabla f(w_{t})\mid\mathcal{F}_{t}]=0$,
\[
\EE\!\big[\|g_{t}\|^{2}\mid\mathcal{F}_{t}\big]
 =\|\nabla f(w_{t})\|^{2}
   +\EE\!\big[\|g_{t}-\nabla f(w_{t})\|^{2}\mid\mathcal{F}_{t}\big]
 \le\|\nabla f(w_{t})\|^{2}+\sigma^{2},
\tag{6}
\]
where the last inequality follows from (A3).

\noindent\textbf{[Step 4]} Plug (6) into (5):
\[
\EE\!\big[f(w_{t+1})\mid\mathcal{F}_{t}\big]
\le f(w_{t})-\eta_{t}\|\nabla f(w_{t})\|^{2}
   +\frac{L\eta_{t}^{2}}{2}\big(\|\nabla f(w_{t})\|^{2}+\sigma^{2}\big).
\tag{7}
\]

\noindent\textbf{[Step 5]} Rearrange (7) to isolate the gradient term:
\[
\big(\eta_{t}-\tfrac{L\eta_{t}^{2}}{2}\big)\|\nabla f(w_{t})\|^{2}
\le f(w_{t})-\EE\!\big[f(w_{t+1})\mid\mathcal{F}_{t}\big]
   +\frac{L\eta_{t}^{2}}{2}\sigma^{2}.
\tag{8}
\]

\noindent\textbf{[Step 6]} Since $\eta_{t}<1/L$ (Assumption A4), we have
$\eta_{t}-\frac{L\eta_{t}^{2}}{2}\ge\frac{\eta_{t}}{2}$.
Thus (8) yields
\[
\frac{\eta_{t}}{2}\,\|\nabla f(w_{t})\|^{2}
\le f(w_{t})-\EE\!\big[f(w_{t+1})\mid\mathcal{F}_{t}\big]
   +\frac{L\eta_{t}^{2}}{2}\sigma^{2}.
\tag{9}
\]

\noindent\textbf{[Step 7]} Take total expectation and sum (9) over $t=0,\dots,T-1$:
\[
\frac{1}{2}\sum_{t=0}^{T-1}\EE[\eta_{t}\|\nabla f(w_{t})\|^{2}]
\le \EE[f(w_{0})]-\EE[f(w_{T})]
   +\frac{L\sigma^{2}}{2}\sum_{t=0}^{T-1}\eta_{t}^{2}.
\tag{10}
\]

\noindent\textbf{[Step 8]} Drop the non‑negative term $\EE[f(w_{T})]\ge f^{\star}$ and denote
$f^{\star}=\inf_{w}f(w)$. Then
\[
\frac{1}{2}\sum_{t=0}^{T-1}\EE[\eta_{t}\|\nabla f(w_{t})\|^{2}]
\le f(w_{0})-f^{\star}
   +\frac{L\sigma^{2}}{2}\sum_{t=0}^{T-1}\eta_{t}^{2}.
\tag{11}
\]

\noindent\textbf{[Step 9]} Divide both sides of (11) by $T\bar\eta$, where $\bar\eta:=\frac{1}{T}\sum_{t=0}^{T-1}\eta_{t}$:
\[
\frac{1}{T}\sum_{t=0}^{T-1}\EE[\|\nabla f(w_{t})\|^{2}]
\le
\frac{2\big(f(w_{0})-f^{\star}\big)}{\bar\eta\,T}
   +\frac{L\sigma^{2}}{\bar\eta\,T}
        \sum_{t=0}^{T-1}\eta_{t}^{2}.
\tag{12}
\]

\noindent\textbf{[Step 10]} Use the deterministic bounds $\eta_{\min}\le\eta_{t}\le\eta_{\max}$ to bound the sum of squares:
\[
\sum_{t=0}^{T-1}\eta_{t}^{2}\le T\,\eta_{\max}^{2},
\qquad
\sum_{t=0}^{T-1}\eta_{t}^{2}\ge T\,\eta_{\min}^{2}.
\]
Plugging the upper bound into (12) gives the inequality stated in Theorem \ref{thm:main}:
\[
\frac{1}{T}\sum_{t=0}^{T-1}\EE[\|\nabla f(w_{t})\|^{2}]
\le
\frac{2\big(f(w_{0})-f^{\star}\big)}{\bar\eta\,T}
  +L\frac{\eta_{\max}^{2}}{\bar\eta}\sigma^{2}.
\]
A symmetric argument using the lower bound yields the term $L\eta_{\min}^{2}\sigma^{2}/\bar\eta$; adding both gives the compact form
\[
\frac{2\big(f(w_{0})-f^{\star}\big)}{\bar\eta\,T}
  +L\big(\eta_{\max}^{2}+\eta_{\min}^{2}\big)\frac{\sigma^{2}}{\bar\eta},
\]
which is exactly (3). ∎

\section*{Rate Derivation (With Constants)}
If we choose a *triangular* schedule with amplitude shrinking as
\[
\eta_{\max}=c_{1}T^{-1/2},\qquad
\eta_{\min}=c_{2}T^{-1/2},
\]
then $\bar\eta=\Theta(T^{-1/2})$ and $\eta_{\max}^{2}+\eta_{\min}^{2}=O(T^{-1})$. Substituting into (3) yields
\[
\frac{1}{T}\sum_{t=0}^{T-1}\EE[\|\nabla f(w_{t})\|^{2}]
=O\!\left(\frac{1}{\sqrt{T}}\right).
\]

\section*{Special Cases}
\begin{itemize}
\item \textbf{Constant step size} ($\eta_{\min}=\eta_{\max}=\eta$):
    \[
    \frac{1}{T}\sum_{t=0}^{T-1}\EE[\|\nabla f(w_{t})\|^{2}]
    \le\frac{2(f(w_{0})-f^{\star})}{\eta T}
      +L\eta\sigma^{2},
    \]
    recovering the classic SGD bound.
\item \textbf{Purely increasing schedule} (no down‑ramp):
    The proof remains valid because only the bounds $\eta_{\min},\eta_{\max}$ are used; the cyclic nature is not essential for the rate, but it may improve practical performance.
\end{itemize}

\end{document}